\documentclass[11pt,a4paper]{article}
\usepackage[T1]{fontenc}
\usepackage[utf8]{inputenc}
\usepackage[polish]{babel}
\usepackage{lmodern}
\usepackage[a4paper,margin=2.7cm]{geometry}
\usepackage{listings}
\usepackage{hyperref}
\title{Dotrenowanie modeli językowych w zadaniu iSTS\\\large Dokumentacja wstępna}
\author{Wiktor Frątczak, \newline
Piotr Hołownia}
\date{\today}
\begin{document}
    \maketitle


    \section{Cel projektu}
    Celem projektu jest zbadanie skuteczności dotrenowanych
modeli językowych w zadaniu interpretowalnego podobieństwa
semantycznego (iSTS) dla par zdań w języku angielskim.
System ma wskazać dopasowania fragmentów zdań, typy relacji
oraz oceny podobieństwa. Wyniki zostaną porównane z
wcześniejszymi rozwiązaniami opisanymi w materiałach
projektowych.


    \section{Zakres zadania}
    Skupimy się na wariancie SysChunks, w którym system sam
wyznacza fragmenty zdań. Planujemy dotrenowanie metodami
LoRA i QLoRA oraz porównanie z pełnym dotrenowaniem
mniejszych modeli, jeśli dostępne zasoby na to pozwolą.

    Ewaluacja obejmie dostępne zbiory SemEval 2016 iSTS z
domen \textit{images}, \textit{headlines} i
\textit{answers-students}. Każdy eksperyment zostanie
uruchomiony pięć razy z różnymi wartościami ziarna
losowości. Porównania zostaną uzupełnione analizą istotności
statystycznej.


    \section{Model i plan eksperymentów}
    Planujemy eksperymenty z modelami Gemma 4 31B oraz
GPT-6.1 SOL. Model SOL wykorzystamy do ewaluacji jako punkt
odniesienia. \textbf{Do potwierdzenia pozostają dokładne
warianty modeli, dostęp do GPU oraz możliwość uruchomienia i
dotrenowania Gemmy.} Dla Gemmy porównamy wynik modelu
bazowego z wynikami po dotrenowaniu metodami LoRA i QLoRA. W
eksperymentach wykorzystamy udostępniony kod projektu
Jacyny.

    Dla Gemmy i wariantu SysChunks planujemy pięć uruchomień
dla każdej z trzech konfiguracji: modelu bazowego, LoRA i
QLoRA, czyli łącznie 15 uruchomień na domenę. Dla baseline'u
planujemy powtórzenie inferencji, a dla adapterów osobne
treningi i ewaluacje. SOL będzie oceniany oddzielnie jako
model referencyjny. \textbf{Do sprawdzenia pozostaje obsługa
seedów inferencji przez wybrany backend i jej wpływ na
powtarzalność.} Zmienność baseline'u będzie zależała od
ustawień generowania. W miarę dostępności zasobów rozważymy
także pełne dotrenowanie mniejszych modeli.


    \section{Przygotowanie danych}
    Dane iSTS zawierają pary zdań i referencyjne anotacje
dopasowań fragmentów. Trzy domeny obejmują nagłówki
wiadomości, opisy obrazów oraz odpowiedzi studentów
porównywane z odpowiedziami wzorcowymi. W dostarczonych
materiałach dostępne są zbiory treningowe i testowe oraz
referencyjne anotacje. Przygotowanie przykładów do
dotrenowania nadzorowanego będzie polegało na określeniu
wejścia obejmującego parę zdań, instrukcję systemową.
Oczekiwana odpowiedź zostanie przygotowana na podstawie
dostępnych anotacji.

    \subsection{Prompt i format odpowiedzi}
    W aktualnych testach wariant zero-shot JSON korzysta z
poniższej instrukcji systemowej zapisanej w
\texttt{prompts/system/zero-shot.json.txt}. Instrukcja
określa typy relacji, skalę ocen i format odpowiedzi. Skrypt
przekazuje surową parę zdań w osobnej wiadomości
użytkownika, a prompt zleca modelowi numerowanie tokenów
rozdzielonych białymi znakami od 1 osobno dla każdego
zdania.

    \begin{lstlisting}[basicstyle=\small\ttfamily,breaklines=true,breakatwhitespace=true,columns=fullflexible,keepspaces=true,frame=single]
Perform Interpretable Semantic Textual Similarity (iSTS) for
the supplied English sentence pair.

Use these relation labels.
- EQUI means the chunks are equivalent in context.
- OPPO means their meanings are opposed in context.
- SPE1 means the chunk in Sentence 1 is more specific.
- SPE2 means the chunk in Sentence 2 is more specific.
- SIMI means the meanings are similar without one of the
relations above.
- REL means the chunks are related without one of the
relations above.
- NOALI means the chunk has no counterpart in the other
sentence.

Use an integer score from 0 to 5. Use 5 for equivalent
meanings,
3 or 4 for very similar or closely related meanings, 1 or 2
for slightly
similar or related meanings, and 0 for unrelated or
unaligned chunks.

Segment each sentence into meaningful chunks, then align
corresponding chunks.

Number the whitespace-separated tokens starting at 1
independently in each
sentence. Use the original token order. List the tokens that
participate in
each alignment. For NOALI, use an empty array for the
missing side and a
score of 0.

Return a single JSON object with only an alignments array.
Each row must
contain source_indices, target_indices, label, and score.
Use integer token
indices. Return plain JSON without Markdown fences,
comments, or explanation.

{
  "alignments": [
    {
      "source_indices": [1, 2],
      "target_indices": [1],
      "label": "SPE1",
      "score": 4
    }
  ]
}
    \end{lstlisting}

    Odpowiedź JSON ma zawierać tablicę \texttt{alignments},
a każda pozycja listy indeksów, etykietę i ocenę. Wybraliśmy
JSON, ponieważ nazwane pola ułatwiają parsowanie i
walidację. Pakiet \texttt{ists\_converter} zawiera konwersję
w obie strony między JSON i formatem \texttt{.wa}, w tym
mapowanie oceny NIL na 0 i odwrotnie dla niedopasowanych
fragmentów. Konwerter pomija oryginalne komentarze anotacji
i generuje je ponownie. Dla niedopasowanych fragmentów
normalizuje zapis oceny 0 do NIL. Testy sprawdzają wybrane
przypadki konwersji, bez potwierdzenia identyczności
wszystkich plików bajt po bajcie. W próbnej ewaluacji
podanej w notatkach odpowiedzi JSON były poprawne składniowo
dla wszystkich 1094 par.

    Walidator sprawdza strukturę pól, indeksy, dozwolone
etykiety i zakres oceny. Dopuszcza niepełne pokrycie słów i
nakładające się dopasowania oraz nie wymusza oceny 0 dla
NOALI, choć wymaga jej prompt. Przejście walidacji nie
potwierdza więc spełnienia wszystkich zasad zadania. Skrypt
ewaluacyjny zapisuje odpowiedź, której nie udało się
przekonwertować, i dla danej pary przekazuje do ewaluacji
pustą listę dopasowań. Błąd żądania lub nieprawidłowa
odpowiedź API mogą przerwać uruchomienie. Skrypt dopuszcza
ręczne poprawki zapisanych odpowiedzi. \textbf{Do ustalenia
pozostaje końcowy protokół obsługi błędnych odpowiedzi, w
tym zasady ręcznych poprawek.}

    Konwerter obsługuje opcjonalne znaczniki FACT i POL, ale
powyższy prompt ich nie wymienia. \textbf{Do ustalenia
pozostaje sposób ich uwzględnienia w instrukcji oraz
zgodność przygotowania danych treningowych z wybranym
formatem odpowiedzi.}

    \subsection{Typy etykiet}
    Podstawowe typy etykiet to EQUI, OPPO, SPE1, SPE2, SIMI,
REL i NOALI. Dodatkowe znaczniki FACT i POL opisują
odpowiednio różnice w faktyczności oraz polaryzacji.
Specyfikacja dopuszcza dopasowania wielu fragmentów do wielu
fragmentów.

    \begin{itemize}
        \item \textbf{EQUI} --- fragmenty są równoważne znaczeniowo w kontekście zdań.
        \item \textbf{OPPO} --- znaczenia fragmentów są przeciwstawne.
        \item \textbf{SPE1} --- fragment pierwszego zdania jest bardziej szczegółowy niż odpowiadający mu fragment drugiego.
        \item \textbf{SPE2} --- fragment drugiego zdania jest bardziej szczegółowy niż odpowiadający mu fragment pierwszego.
        \item \textbf{SIMI} --- fragmenty mają podobne znaczenie, ale relacja nie należy do powyższych typów.
        \item \textbf{REL} --- fragmenty są powiązane znaczeniowo, lecz nie spełniają warunków relacji SIMI ani pozostałych wymienionych typów.
        \item \textbf{NOALI} --- fragment nie ma odpowiednika w drugim zdaniu.
    \end{itemize}

    \subsection{Charakterystyka zbiorów}
    Liczebności poniżej uzyskano przez zliczenie par w
lokalnych plikach zdań i zweryfikowano względem liczby
bloków anotacji. Anotacja to wiersz opisujący dopasowanie
fragmentów lub jego brak, a jedna para zdań może zawierać
wiele anotacji.

    \begin{center}
        \begin{tabular}{lrrrr}
            \hline
            Domena           & Pary train & Pary test & Anotacje train & Anotacje test \\
            \hline
            Headlines        & 756        & 375       & 3979           & 2040          \\
            Images           & 750        & 375       & 4805           & 2462          \\
            Answers-students & 330        & 344       & 1848           & 1840          \\
            \hline
        \end{tabular}
    \end{center}

    \textit{Headlines} zawiera pary nagłówków wiadomości, \textit{images} --- podpisów obrazów Flickr, a \textit{answers-students} --- odpowiedzi studentów i odpowiedzi wzorcowych z dziedziny podstaw elektryczności. W każdym zbiorze dominują anotacje NOALI i EQUI. Przykładowo w danych treningowych \textit{headlines} występują odpowiednio 1792 i 1323 takie anotacje, a w \textit{images} 2892 i 1029. Znaczniki FACT i POL występują lokalnie w \textit{headlines}: odpowiednio 29 i 4 razy w train oraz 0 i 2 razy w test. Liczby dotyczą wierszy anotacji, a nie niezależnych par zdań.

    \subsection{Rozkład anotacji i jakość danych}
    Tabela przedstawia liczbę wierszy anotacji dla każdego
podstawowego typu relacji w dostarczonych danych
treningowych. Liczby obejmują również wiersze z dodatkowymi
znacznikami FACT i POL.

    \begin{center}
        \begin{tabular}{lrrrrrrr}
            \hline
            Domena           & NOALI & EQUI & SIMI & SPE1 & SPE2 & REL & OPPO \\
            \hline
            Headlines        & 1792  & 1323 & 324  & 200  & 193  & 128 & 19   \\
            Images           & 2892  & 1029 & 340  & 245  & 229  & 67  & 3    \\
            Answers-students & 930   & 595  & 73   & 61   & 74   & 76  & 39   \\
            \hline
        \end{tabular}
    \end{center}

    Rozkłady są niezrównoważone: NOALI i EQUI występują
najczęściej, a OPPO jest rzadkie, szczególnie w zbiorze
\textit{images}. W danych testowych występuje odpowiednio
13, 1 i 49 anotacji OPPO w domenach \textit{headlines},
\textit{images} i \textit{answers-students}. Liczebności
rzadkich relacji mogą utrudnić ocenę jakości dla
poszczególnych typów. Średnia liczba wierszy anotacji na
parę treningową wynosi odpowiednio 5,26, 6,41 i 5,60
\textbf{TO BE VERIFIED}.

    W testowym \textit{answers-students} para o
identyfikatorze 252 ma pustą sekcję dopasowań. W danych
występują także różne zapisy oceny dla relacji NOALI, w tym
NIL i 0. \textbf{Przed przygotowaniem przykładów należy
sprawdzić interpretację tych przypadków w ewaluatorze i
ustalić ich obsługę.} Dalsza analiza obejmie długości
tokenizowanych wejść i oczekiwanych odpowiedzi oraz zgodność
anotacji z wymaganym formatem.

    Dostarczony podział obejmuje train i test. W obu
zbiorach \textit{headlines} znaleziono cztery identyczne
uporządkowane pary zdań w porównaniu bez normalizacji
tekstu. \textbf{Do ustalenia pozostaje sposób obsługi tych
par, przy zachowaniu oficjalnego zbioru testowego.}
\textbf{Do ustalenia pozostają sposób wydzielenia zbioru
walidacyjnego i jego proporcja.} Zliczenia par treningowych
dla \textit{headlines} i \textit{answers-students} różnią
się od odpowiednio 750 i 333 podanych w pracy Ścisla.
\textbf{Przyczyna rozbieżności pozostaje niepewna i wymaga
sprawdzenia.}


    \section{Dotrenowanie i walidacja odpowiedzi}
    Zastosujemy nadzorowane dotrenowanie, w którym model
otrzymuje przygotowane wejście i odpowiadającą mu anotację.
LoRA i QLoRA będą trenowane jako osobne adaptery. Przed
eksperymentami ustalimy m.in. współczynnik uczenia, liczbę
epok, rozmiar partii, długość sekwencji oraz rangę,
skalowanie i dropout adapterów. Hiperparametry dobierzemy na
podstawie dokumentacji wybranego wariantu Gemmy oraz
dostępnych zasobów GPU. Wartości dobierzemy na zbiorze
walidacyjnym, zapisując konfiguracje i uzasadnienie wyboru.

    Podczas pomiaru baseline'u sprawdzimy, czy odpowiedzi
zachowują wymagany format, czy model nie dodaje dodatkowych
wyjaśnień lub samopoprawek oraz czy odpowiedzi nie są
ucięte. \textbf{Do ustalenia pozostaje sposób walidacji i
obsługi niepoprawnych odpowiedzi przed właściwym
porównaniem.} Odsetek błędnych odpowiedzi będzie
raportowany.


    \section{Dane i ewaluacja}
% Uzyjemy skryptwo distarczonych z datasetami, oraz promptów. Bedziemy testowac nasze modele z ich promptami do przetestowania baselinu, aby zobaczyc jak modele sobie radza

    Skrypt \texttt{evaluate\_ists.py} wykorzystuje dwa
dostarczone warianty ewaluatora F1. Raportuje F1
dopasowania, F1 typu i F1 oceny z wariantu no penalty oraz
F1 łączące typ i ocenę dla obu wariantów. Skrypt
\texttt{wellformed.pl} służy do kontroli formatu. Notebook
Colab uruchamia inferencję przez vLLM i skrypt ewaluacyjny;
jego ustawienia demonstracyjne obejmują model Qwen2.5 1.5B
oraz pięć par Images. Konfiguracja ta służy sprawdzeniu
działania pipeline'u. Notebook i skrypt ewaluacyjny nie
realizują treningu LoRA ani QLoRA, a ich obecne ustawienia
nie potwierdzają uruchomienia Gemmy.

    \subsection{Wstępna ewaluacja SOL 6.1}
    W przekazanych wynikach próbnej ewaluacji modelu GPT-6.1
SOL porównano prompty zero-shot, one-shot i eight-shot. W
przekazanej tabeli dla 1094 par wynik zbiorczy F1 typu i
oceny w wariancie no penalty wyniósł 0,5828 dla
\texttt{zero-shot.ists}, 0,5977 dla
\texttt{zero-shot.jacyna-json}, 0,6036 dla
\texttt{one-shot.jacyna-json} oraz 0,6372 dla
\texttt{eight-shot.ists}. Dla tego ostatniego promptu wyniki
w domenach Images i Headlines wyniosły odpowiednio 0,6381 i
0,6258.

    W czterech sprawdzonych wariantach JSON nie odnotowano
błędnych składniowo odpowiedzi ani odrzuconych wyników. W
wariantach \texttt{zero-shot.ists} i \texttt{one-shot.ists}
odrzucono po 63 z 1094 odpowiedzi. Wyniki te stanowią
wstępny punkt odniesienia dla dalszych eksperymentów.
\textbf{Do potwierdzenia pozostają identyfikator i
ustawienia modelu, wariant segmentacji, sposób obliczenia
wyniku zbiorczego oraz obsługa odrzuconych odpowiedzi w tej
ewaluacji.} Skrypt zapisuje wyniki osobno dla domen, a
sposób wyliczenia przekazanej tabeli zbiorczej wymaga
weryfikacji. Pierwsza tabela wyników domenowych nie została
podpisana, dlatego nie przypisujemy jej domeny.


    \section{Plan analizy statystycznej}
    Dla LoRA i QLoRA planujemy po pięć treningów i ewaluacji
z różnymi ziarnami losowości. Dla modelu bazowego planujemy
powtórzyć inferencję bez dotrenowania, przy czym wpływ
ziarna zależy od backendu i ustawień generowania. Dla
każdego uruchomienia zapiszemy wynik F1, konfigurację oraz
predykcje dla poszczególnych par zdań. Następnie
przedstawimy wszystkie wyniki, ich średnią i odchylenie
standardowe. Pozwoli to opisać zmienność wyników pomiędzy
uruchomieniami.

    Przy porównywaniu eksperymentów planujemy podawać
różnicę wyników wraz z analizą jej istotności statystycznej.
Rozważamy także sparowany \textbf{?????????} bootstrap par
zdań testowych: w kolejnych próbach te same pary będą
losowane ze zwracaniem dla obu porównywanych systemów, a
zbiorcze F1 zostanie przeliczone dla każdej próby. Taka
analiza dotyczy zmienności związanej z doborem przykładów
testowych. Oddzielnie trzeba uwzględnić zmienność między
uruchomieniami treningu.

    \textbf{Do ustalenia pozostają jednostka analizy, metoda uwzględnienia obu źródeł zmienności, próg istotności oraz sposób obsługi wielokrotnych porównań. - NOTE} Wyniki planujemy przedstawić wraz z wielkością zaobserwowanej różnicy i miarą niepewności.


    \section{Kwestie do ustalenia}
    \textbf{Do ustalenia pozostają dostęp do GPU, dokładny wariant modelu, zakres pełnego dotrenowania, format danych treningowych i predykcji, sposób obsługi niepoprawnych odpowiedzi, metoda analizy statystycznej oraz protokół porównania z wynikami referencyjnymi.}

    \begin{thebibliography}{9}
        \bibitem{semeval} SemEval 2016 Task 2, \url{https://alt.qcri.org/semeval2016/task2/}.
        \bibitem{kopyt} M. Kopyt, \textit{Określanie interpretowalnego podobieństwa semantycznego dla języka angielskiego}, praca magisterska, rozdział 4.
        \bibitem{jacyna} Z. Jacyna, P. Pusiewicz, M. Rak, \textit{Dokumentacja końcowa projektu dotyczącego podobieństwa semantycznego}.
        \bibitem{gemma} Google, dokumentacja Gemma, \url{https://ai.google.dev/gemma/docs/core}.
        \bibitem{scisel} A. Ścisel, praca magisterska dostarczona w materiałach projektu.
    \end{thebibliography}
\end{document}
